\section{Introduction}
\label{sec:introduction}

Urban-scale wind and pollutant prediction depends critically on
inflow forcing that is rarely known with certainty.  Inflow direction,
speed, and bulk pressure gradient drift on time scales comparable to
those of the flow itself, so any parameter recovered from
measurements must in principle be treated as time-varying.  At the
same time, large-eddy simulations (\LES{}) of urban flows are too
expensive to admit a fully nonlinear Bayesian inversion: ensemble
sizes are limited to a few tens of members, the assimilation window
has to be kept short, and each forward run is a non-differentiable
black box.

A pragmatic answer to this combination of constraints is the
ensemble smoother with multiple data
assimilation~(\ESMDA{})~\cite{emerick2013esmda}, an ensemble Kalman
method~\cite{evensen2003ensemble,evensen2009book} that approximates a
single Bayesian update by a sequence of tempered Kalman steps.  As
an iterative ensemble smoother~\cite{stordal2018itesmda}, \ESMDA{}
does not require gradients and remains affordable for ensembles of a
few tens of members.  Applied to a long simulation horizon, however,
\ESMDA{} on a single window quickly becomes intractable -- the
parameter dimension grows linearly with the number of representation
points and the cost of each forward sweep grows with the simulated
time.  A natural remedy is to split the horizon into short windows
and assimilate sequentially, propagating posterior information
forward through a generative model of the parameter trajectory.

This report describes such a workflow.  We combine \ESMDA{}, applied
to a piecewise representation of the inflow parameters, with a
parameter time-series model that draws the cold-start prior of the
first window and propagates the posterior of one window into the
prior of the next.  Continuity of the recovered trajectory across
window boundaries is enforced by two complementary mechanisms: the
time-series model itself, which can produce $C^1$-continuous
predictions per ensemble member, and an option in the smoother to
\emph{pin} the window-initial parameter so that the Kalman update
does not adjust it.  The construction extends the rolling-window
\ESMDA{} of Evensen~\cite{evensen2024} to a software design in
which the prior generator and the smoother are decoupled: four
generative models -- a Gaussian-process prior with linear-trend
extrapolation~\cite{rasmussen2006gp}, a stationary
auto-regressive AR(1) prior, an Ornstein--Uhlenbeck
prior~\cite{uhlenbeck1930ou}, and a critically-damped AR(2)
relaxation following Evensen~\cite{evensen2024} -- can be used
interchangeably without touching the assimilation logic.

The implementation lives in the \pyurbanair{} repository.  Three
forward models are supported behind a uniform Python interface: a
custom lattice-Boltzmann solver~(\pylbm{}), the urban-flow \LES{}
code uDALES~(\pyudales{})~\cite{suter2018udales}, and
PALM~(\pypalm{})~\cite{maronga2020palm}.  The data-assimilation
pipeline is solver-agnostic and interacts with each solver only
through ensemble forward and observation operators.  This report
focuses on the workflow driven by
\rolloutscript{}, which is
the script that we ultimately want to publish on.

\paragraph{Outline.}
Section~\ref{sec:problem} states the inverse problem and the
identical-twin experimental setting.  Section~\ref{sec:forward}
describes the forward models and the common interface they expose.
Section~\ref{sec:da} introduces the data-assimilation strategy in
three subsections: single-window \ESMDA{} for time-varying
parameters, prior parameter generation, and the multi-window
rollout.  Section~\ref{sec:results} reports the experimental setup
and the diagnostics that the rollout produces, and
Section~\ref{sec:conclusion} concludes with planned extensions.
